\section{Threats to Validity}
\label{sec:threats}

We discuss threats to validity following the taxonomy of Wohlin et al.~\cite{wohlin2012experimentation}.

\subsection{Internal Validity}

\textbf{Model version drift.}
The \texttt{unsloth/Qwen2.5-7B-Instruct} checkpoint on HuggingFace
may be silently updated between setup and the full run, affecting
reproducibility. We mitigated this by pinning the exact model
revision at download time and logging all inference configuration
in \texttt{results/full\_api\_log.txt}.

\textbf{Non-zero temperature.}
Qwen2.5-7B-Instruct requires \texttt{temperature~>~0}; we used the
minimum permitted value of 0.01. To ensure deterministic decoding
we set \texttt{do\_sample=False} (greedy), which suppresses the
effect of temperature on output selection.

\textbf{Unspecified GPT-4o version in reference paper.}
E02~\cite{acharya2025bugquality} reports threshold values
(CTQRS=75\%, SBERT=0.82, ROUGE-1=0.47) attributed to
``ChatGPT-4o'' without specifying the exact model version string
(e.g., \texttt{gpt-4o-2024-08-06}). Notably, reference [53] within
E02 cites ``ChatGPT (Version 3.5)'' while Figure~4 reports
``ChatGPT 4o'', indicating an internal inconsistency in the source
paper. We therefore treat the published thresholds as historical
reference points rather than guaranteed reproducible values tied
to a fixed model version.

\subsection{External Validity}

\textbf{Single domain.}
Our test set is drawn exclusively from the Bugzilla/Mozilla domain.
Results may not generalize to other bug tracking systems or
application domains. We scoped our conclusions to the Bugzilla
domain and note cross-domain generalization as future work
(Section~\ref{sec:conclusion}).

\textbf{Limited cross-domain data.}
Our study does not include a cross-domain evaluation (originally
planned as RQ2 using 37 Mojira/Minecraft pairs from
E03~\cite{akyol2026improbr}). On closer inspection of the
full text of E03, those 37 pairs constitute a curated internal
set used for the ImproBR/ChatBR comparison and are not
available as a reusable public corpus. We therefore excluded RQ2
from the final scope and flag cross-domain validation as an open
research direction.

\subsection{Construct Validity}

\textbf{Non-overlapping test splits.}
Our test set (N=262) was produced by a 80/10/10 split with
\texttt{random\_state=42}; E02 used N=3,903 after an additional
SBERT-based deduplication step with an undisclosed seed,
yielding N=391 test samples. The two test sets are unlikely to be
identical, so published thresholds represent reference baselines
rather than results on the same data.

\textbf{Absence of paired baseline.}
We cannot perform paired significance testing between Qwen and
GPT-4o because per-instance GPT-4o outputs are not publicly
available. This limitation is acknowledged by prior
work~\cite{choi2025agentreport}, which states that
``paired significance testing not feasible due to unavailability
of per-instance baseline results.'' We adopt one-sample
Wilcoxon tests against the published threshold medians as
the methodologically sound alternative.

\textbf{CTQRS scorer discrepancy.}
The GindeLab scorer~\cite{acharya2025bugquality} hard-codes
\texttt{max\_possible=16}, while E02 reports a maximum of 17 in
the paper text. Computing the true maximum from the 13
sub-check weights in the scorer source code yields 18 (RM1--4:
$4\times1$; RR1/RR2/RR4/RR5: $4\times1$; RR3: $1\times2$;
RA1--4: $4\times2$). We used \texttt{max\_possible=18} as the
only value derivable from first principles without ambiguity,
ensuring CTQRS percentages remain bounded in $[0, 100]$.
We also corrected word-list noise in the RA sub-checks
(tokens shorter than 3 characters or containing whitespace
were removed) and applied whole-word boundary matching
to prevent spurious matches; details are in
\texttt{scripts/ctqrs\_scorer\_v2\_fixed.py}.

\textbf{Metric sensitivity to style.}
SBERT (\texttt{paraphrase-MiniLM-L6-v2}) encodes entire
passages into fixed-size vectors and is sensitive to differences
in length and formatting style, not only semantic content.
Our ablation (Section~\ref{sec:discussion}) confirms that
Qwen outputs are approximately 1.82$\times$ longer than
ground-truth reports and use markdown headers absent
from the plain-text Bugzilla reference reports, which
depresses cosine similarity independently of content quality.
CTQRS and ROUGE-1 provide complementary perspectives
that are less affected by this style gap.

\subsection{Conclusion Validity}

\textbf{Multiple testing.}
RQ1 involves three separate Wilcoxon tests (CTQRS, SBERT,
ROUGE-1), inflating the familywise Type~I error rate.
We applied Bonferroni correction ($\alpha_{\text{adj}} =
0.05/3 \approx 0.017$) and report both raw and
corrected $p$-values; the conclusions are unchanged
after correction.

\textbf{Effect size reporting.}
$p$-values alone do not quantify practical importance.
We therefore report Cliff's $\delta$ alongside every
Wilcoxon test~\cite{arcuri2011practical}, interpreted
using the conventional thresholds ($|\delta|<0.147$:
negligible; $<0.33$: small; $<0.474$: medium; otherwise
large).
